% =========================================================================
% Section I -- Introduction.
% Claims are deliberately scoped to the cited formulations and to the
% controlled instance family. Insert numerical findings only after the final
% experimental protocol has been frozen and executed.
% =========================================================================

\section{Introduction}
\label{sec:intro}

\begin{figure}[t]
  \centering
  \includegraphics[width=\columnwidth]{fig_layout}
  \caption{Fixed $20\times19$ layout used throughout. Five inbound and five
  outbound stations have single-file, depth-three waiting lines. Sixteen
  dedicated charging bays occupy four aisle rows. The dashed arrow illustrates
  one pickup--delivery pair.}
  \label{fig:layout}
\end{figure}

Fleets of autonomous mobile robots (AMRs) increasingly perform internal
transfers in factories and terminals. In the fixed layout studied here
(Fig.~\ref{fig:layout}), jobs
arriving at inbound stations must be consolidated on capacitated robots and
delivered to outbound stations across a shared floor. The scheduler selects
the carrier and order of each pickup and delivery. Its decisions are evaluated
in an execution setting in which robots require collision-free trajectories
and each station serves at most one robot at a time and provides only finite
waiting space.

Learned constructive schedulers encode partial schedules with graph or
attention models and have progressed from operation--machine decisions to
transport-aware operation--machine--vehicle choices
\cite{zhang2020l2d,song2023fjsp,moon2024hgs,yuan2026gin,matrans2026}. These
methods commonly represent transport through fixed pairwise travel times. In
our pickup-and-delivery setting, that abstraction does not reproduce delays
caused by route reservations or non-preemptive contention at shared inbound
and outbound stations.

Assignment and routing have also been integrated directly, through routed-cost
task assignment and reinforcement-learning frameworks coupled with
conflict-free path planners
\cite{chen2021integrated,li2024hierarchical,wang2026collaborative}. These studies 
establish the value of routing-aware assignment, but their principal questions 
concern online allocation or per-instance optimization rather than how route conflicts 
and contention at exclusive stations with finite waiting space affect evaluator fidelity 
over a common set of complete offline schedules. This paper addresses that complementary 
question by quantifying how an idealized
travel-time model distorts both the predicted makespan and the ranking of
complete schedules when robot capacity, non-zero station service times, finite
waiting space, and collision-aware execution interact.

We study this question through a paired comparison. For the same assignment
and operation sequence, an idealized decode charges free-space travel and
starts service on arrival, whereas the fixed collision-aware executor plans
space--time paths and enforces exclusive service. By evaluating the same
schedule under both models, the execution penalty isolates
makespan-estimation error from differences in schedule construction. More
importantly, we compare the orderings induced by the two evaluators: if they
disagree, the idealized evaluator is not merely miscalibrated but can prefer a
schedule that executes worse. The executor jointly captures route-reservation
delays and contention at exclusive stations; we evaluate their combined effect
on schedule value and ranking rather than attribute the gap to individual
delay sources.

The same distinction guides the learning method. We train a constructive
policy with PPO using a terminal reward derived from makespan returned by the
fixed collision-aware executor. During both training and inference, however,
partial schedules are advanced by a calibrated event-based surrogate rather
than by the executor. The policy therefore observes surrogate-projected station
availability, queue occupancy, and outstanding committed workload, together
with projected local robot density and an estimated wait for each candidate
action. The surrogate's affine travel-time correction and
queue-depth-conditioned waiting penalties are fitted offline to
operation-level traces generated by the fixed collision-aware executor. This
calibration is designed to reduce systematic drift in the intermediate robot
and station clocks while avoiding space--time routing at every construction
step. Once a schedule is complete, the fixed executor evaluates its makespan;
during training, this executor-evaluated cost supplies the terminal reward.
Thus, execution grounding refers to the optimization target, whereas
intermediate state construction remains surrogate-based.

On the 60-job test set, the idealized and executed rankings have mean
$\tau_b=0.334\pm0.026$ and invert $33.0\pm1.3\%$ of comparable pairs; choosing
by the idealized evaluator incurs $29.8\pm3.7\%$ execution regret within the
comparison set. Executor-grounded training is the dominant factor in the
ablation, while the benefit of exposing projected congestion grows from the
40- to the 100-job setting. The three-seed policy with eight executor-scored
rollouts reduces makespan by $5.9\%$ relative to the test-set best fixed rule
and by $3.6\%$ relative to the genetic baseline at 60 jobs.

This study asks when an idealized travel-time model fails to preserve the
ranking of complete schedules, and whether executor-grounded learning with
surrogate-projected congestion features can reduce the resulting
schedule-selection error.
Specifically, this paper makes three contributions:

\begin{itemize}

\item A paired evaluation of complete offline capacitated AMR schedules under
an idealized transport decoder and a fixed collision-aware executor. The
evaluation quantifies value error through the execution penalty and ranking
error through Kendall's $\tau_b$ and the directly measured ranking inversion
rate.

\item A controlled workload-scale experiment that fixes the fleet and layout,
varies the number of jobs, and maps the non-monotone regimes in which the two
evaluators rank complete schedules differently. Value error, ranking
inversion, selection regret, and routing-failure rate characterize the
associated conditions.

\item An executor-grounded constructive policy trained with PPO using a
terminal reward derived from makespan returned by the fixed collision-aware
executor. During both training and inference, a calibrated event-based
surrogate advances partial schedules and supplies projected station states,
local robot density, and action-specific waiting estimates. A $2\times2$
factorial evaluation crosses this surrogate-projected congestion representation
with the source of the training objective at matched computational budgets.

\end{itemize}

Section~\ref{sec:related} reviews learned scheduling and integrated task and
path planning. Section~\ref{sec:problem} defines the scheduling problem, the
two evaluators, and the value and ranking measures. Section~\ref{sec:architecture}
presents the calibrated surrogate, congestion-aware policy architecture, and
PPO training procedure. Section~\ref{sec:experiments} reports the schedule-
ranking analysis, workload-scale study, policy evaluation, and ablations.
Section~\ref{sec:conclusion} concludes the paper.
